\section{Introduction}

Legal LLM evaluation often collapses a model's answer into a single correctness
score. That misses a central feature of legal reasoning: different answers may
reach the same bottom line through materially different gates, exceptions,
factual characterizations, or counterarguments. A benchmark that only grades
individual responses can therefore obscure whether models understand the
structure of the doctrine or merely produce plausible conclusions.

This paper proposes a source-grounded pipeline for evaluating legal reasoning
paths. The pipeline starts from a source case and automatically constructs the
benchmark artifacts needed to test models: a legal question, gold answer,
controlled variations, a rubric, response clusters, centroid scores, projected
member scores, model rankings, and escalation packets. Statute of Frauds is the
first calibration domain because it supplies clear doctrinal gates and boundary
facts for testing the method. The system is designed, however, to be
doctrine-general: the live path asks Frank and Karthic to infer doctrine and
rubric structure from the source case instead of hard-coding Statute of Frauds
labels.

The goal of the reported study is pre-expert-review internal validation. The manuscript asks
whether the pipeline is coherent, reproducible, robust enough to inspect, and
technically defensible using only automated artifacts and internal review. The
claim at this stage is deliberately bounded but complete: the pipeline can run
end to end on a source case, produce auditable artifacts for every stage, cluster
legal reasoning signatures at controlled scale, apply rubric-based LLM judging,
project centroid scores, generate model rankings, and surface failures through
explicit escalation packets.

The validation is organized around the questions a research team would ask while
reviewing the run bundle. Did Frank produce the expected legal framing, neutral
question, gold answer, and variations? Did Karthic convert that packet into a
comprehensive rubric rather than a generic checklist? Did the response models
answer naturally? Did Dasha cluster the resulting answers by legal theory? Did
the judge apply the rubric faithfully enough to make projected rankings
auditable? The manuscript treats those stage-level questions as part of the
method, not as informal commentary outside the paper.
